\documentclass[
  12pt,
  openright,
  oneside,
  a4paper,
  brazil
]{abntex2}

\usepackage[T1]{fontenc}
\usepackage[utf8]{inputenc}
\usepackage{newtxtext,newtxmath}
\usepackage{microtype}
\usepackage{indentfirst}
\usepackage{graphicx}
\usepackage[alf]{abntex2cite}
\usepackage[hidelinks]{hyperref}

\setlrmarginsandblock{3cm}{2cm}{*}
\setulmarginsandblock{3cm}{2cm}{*}
\checkandfixthelayout
\setlength{\parindent}{1.25cm}
\setlength{\parskip}{0cm}

\titulo{Detecção e Explicação de Risco de Fraude em Contexto Pix com Machine Learning, SHAP e RAG}
\autor{Letícia Alves \\ Lucas Nogueira}
\local{\texttt{\{\{PREENCHER\}\}}}
\data{2026}
\instituicao{%
  \texttt{\{\{PREENCHER\}\}} \\
  Curso de Ciência da Computação}
\tipotrabalho{Trabalho de Conclusão de Curso}
\preambulo{Trabalho de Conclusão de Curso apresentado como requisito parcial para obtenção do grau de Bacharel em Ciência da Computação.}
\orientador{\texttt{\{\{PREENCHER\}\}}}

\begin{document}
\frenchspacing
\OnehalfSpacing

\pretextual
\imprimircapa
\imprimirfolhaderosto*

\pdfbookmark[0]{\contentsname}{toc}
\tableofcontents*
\cleardoublepage

\textual
\chapter{Introdução}

\section{Contextualização}

O Pix consolidou-se como infraestrutura relevante do sistema de pagamentos brasileiro ao permitir transferências e pagamentos instantâneos em disponibilidade contínua. Em 2024, foi o instrumento de pagamento com maior crescimento em quantidade de transações no país, com aumento de 52\%, e respondeu por 47\% das transações de pagamento sem uso de espécie no último trimestre do ano \cite{bcb_pagamentos_2024}. A ampliação desse uso torna necessária a identificação célere de comportamentos de risco, sem que operações legítimas sejam indevidamente interrompidas.

A detecção de fraude financeira constitui um problema de classificação particularmente difícil. Os casos confirmados são raros em comparação às operações legítimas, os padrões podem mudar ao longo do tempo e a confirmação do rótulo pode ocorrer depois da transação. Nessas condições, uma solução pode apresentar alta acurácia global e, ainda assim, falhar na classe de maior interesse \cite{abdallah2016survey,dalpozzolo2018realistic}. Métricas orientadas à classe positiva, como precisão, revocação, F1 e área sob a curva precisão--revocação, são, portanto, mais informativas do que a acurácia isolada \cite{saito2015precision}.

O desempenho preditivo, contudo, não é suficiente em um domínio no qual decisões automatizadas precisam ser examinadas por pessoas. Um escore elevado, desacompanhado dos fatores que o influenciaram, oferece pouco apoio à análise e pode induzir confiança excessiva. O método SHAP permite atribuir a cada variável uma contribuição local para a previsão \cite{lundberg2017unified}. Essa explicação matemática descreve o comportamento do classificador, mas não contém, por si só, contexto normativo ou operacional sobre o Pix.

A arquitetura de \emph{Retrieval-Augmented Generation} (RAG) complementa modelos de linguagem com trechos recuperados de uma coleção documental \cite{lewis2020rag}. Neste trabalho, a recuperação foi projetada como complemento, e não substituta, da explicação SHAP: o modelo estima o risco; SHAP identifica os fatores que influenciaram a saída; e o RAG recupera fontes normativas, operacionais e setoriais relacionadas ao contexto brasileiro. A geração textual, por sua vez, deve permanecer condicionada a essas evidências e declarar insuficiência de suporte quando necessário.

O Banco Central instituiu o Mecanismo Especial de Devolução por meio da Resolução BCB nº 103/2021, posteriormente incorporada ao Regulamento do Pix \cite{bcb_resolucao103_2021,bcb_resolucao1_2020}. Esse histórico evidencia uma exigência para a base documental: preservar a identificação, a versão e a vigência das fontes. A mistura de regras de períodos distintos pode resultar em uma resposta linguisticamente plausível, porém documentalmente incorreta.

\section{Problema de pesquisa}

Modelos de aprendizado de máquina podem reconhecer padrões complexos em dados transacionais, mas sua saída costuma ser técnica e insuficiente para justificar o resultado em linguagem acessível. Em sentido oposto, um modelo de linguagem pode produzir uma narrativa convincente sem fidelidade ao classificador ou à fonte recuperada. O problema consiste em integrar previsão, explicabilidade e recuperação documental sem transformar correlação em causalidade, sem inventar o significado de atributos anonimizados e sem ocultar as limitações do conjunto de dados.

O estudo utiliza o conjunto IEEE-CIS Fraud Detection, composto por transações de comércio eletrônico fornecidas pela Vesta \cite{ieee_cis_2019}. O conjunto não contém transações Pix e parte de seus atributos é anonimizada. Consequentemente, ele serve como base para uma prova de conceito de classificação de risco e integração técnica, mas não permite afirmar que o desempenho observado represente uma operação Pix real.

\subsection{Questão de pesquisa}

\textbf{Como integrar um modelo de detecção de fraude, explicações locais por SHAP e recuperação de documentos por RAG para apoiar explicações rastreáveis em português, preservando a fidelidade ao modelo e explicitando as limitações de uma prova de conceito baseada no IEEE-CIS?}

\section{Objetivos}

\subsection{Objetivo geral}

Desenvolver e analisar um protótipo híbrido que combine aprendizado de máquina, SHAP e RAG para classificar risco de fraude em dados transacionais e apoiar a produção de explicações em português fundamentadas em evidências do modelo e em documentos verificáveis relacionados ao contexto Pix brasileiro.

\subsection{Objetivos específicos}

\begin{enumerate}
  \item preparar e caracterizar o conjunto IEEE-CIS, documentando desbalanceamento, valores ausentes, anonimização e diferenças em relação ao domínio Pix;
  \item treinar e comparar XGBoost, Random Forest e uma linha de base por meio de métricas adequadas à classe rara;
  \item aplicar SHAP para identificar, global e localmente, as variáveis que influenciam as previsões;
  \item manter um registro de atributos semânticos que impeça a tradução indevida de colunas anônimas para conceitos Pix;
  \item organizar uma base documental versionada com normas, materiais oficiais e fontes setoriais selecionadas;
  \item implementar a recuperação semântica e preparar a geração de explicações com indicação das fontes utilizadas;
  \item definir um protocolo para avaliar a qualidade da recuperação, a coerência da resposta com SHAP, o suporte documental e a clareza para usuários;
  \item disponibilizar código, artefatos permitidos, decisões metodológicas e instruções de reprodução, respeitando limites de licença e privacidade dos dados.
\end{enumerate}

\section{Justificativa}

A Pesquisa Febraban de Tecnologia Bancária 2024 informa que 73\% dos 22 bancos participantes da questão declararam utilizar inteligência artificial na detecção de fraude e lavagem de dinheiro. O relatório também registra estratégias de detecção e resposta entre as prioridades de cibersegurança. Esses resultados não mensuram especificamente fraude no Pix nem representam todo o sistema financeiro, mas evidenciam a pertinência prática de técnicas de inteligência artificial no setor \cite{febraban2024}.

No plano acadêmico, a relevância da proposta está na integração controlada de três tipos de evidência. O escore preditivo indica quanto a observação se aproxima da classe aprendida pelo modelo; SHAP descreve quais entradas influenciaram essa saída; e o RAG acrescenta contexto documental rastreável. A separação explícita dessas evidências reduz o risco de apresentar uma narrativa do modelo de linguagem como se fosse previsão, causalidade ou conteúdo normativo.

O trabalho também discute limitações frequentemente omitidas em protótipos, entre elas a mudança de distribuição, o vazamento de dados durante a preparação, a seleção inadequada de métricas, a temporalidade das normas e o risco de atribuir significado a variáveis anonimizadas. Esse registro evita caracterizar uma demonstração acadêmica como sistema pronto para decisões financeiras em produção.

\section{Delimitação do trabalho}

O escopo compreende uma prova de conceito executada sobre o IEEE-CIS, com modelos tabulares, explicações SHAP, índice vetorial local e interface demonstrativa. Não fazem parte do escopo:

\begin{itemize}
  \item implantação em uma instituição financeira;
  \item processamento de dados pessoais reais de clientes Pix;
  \item decisão automática de bloqueio, acusação ou devolução;
  \item validação jurídica da elegibilidade ao Mecanismo Especial de Devolução;
  \item comprovação de generalização para transações Pix reais;
  \item comparação exaustiva de todos os modelos de detecção de fraude.
\end{itemize}

Os exemplos destinados à interface devem ser anonimizados ou sintéticos. Os documentos regulatórios precisam ser recuperados com metadados de origem, versão e vigência. A saída é tratada como apoio explicativo sujeito à revisão humana. Enquanto não houver cliente de modelo de linguagem configurado, o protótipo deve declarar a indisponibilidade da geração textual, em vez de simular uma explicação.

\section{Contribuições do trabalho}

As contribuições pretendidas são:

\begin{enumerate}
  \item um pipeline reproduzível de preparação, classificação e explicação local;
  \item um esquema de evidências que conecte SHAP à recuperação documental sem inventar semântica para atributos anônimos;
  \item um corpus documental rastreável, com metadados e controle de integridade;
  \item um protocolo de avaliação separado para desempenho preditivo, recuperação e fidelidade das explicações;
  \item uma interface demonstrativa capaz de expor resultados e indisponibilidades de cada camada;
  \item uma discussão transparente sobre a transferência limitada do IEEE-CIS para o contexto Pix.
\end{enumerate}

\section{Organização do trabalho}

Além desta introdução, o Capítulo 2 apresenta os fundamentos de fraude financeira, Pix, aprendizado de máquina, explicabilidade e RAG. O Capítulo 3 descreve os dados, a preparação, os modelos, o protocolo experimental, o corpus e a integração das três camadas. O Capítulo 4 reúne os resultados preditivos e os resultados disponíveis de explicabilidade e recuperação, com suas restrições. Por fim, o Capítulo 5 sintetiza as conclusões, confronta os objetivos com o que foi efetivamente desenvolvido, registra as limitações e indica trabalhos futuros.

\chapter{Revisão de Literatura}

\section{Pix, pagamentos instantâneos e risco de fraude}

O Pix é uma infraestrutura de pagamentos instantâneos operada no âmbito do Sistema de Pagamentos Brasileiro. Sua disponibilidade contínua e ampla adoção tornam relevantes mecanismos capazes de reconhecer sinais de risco com baixa latência. Em 2024, o instrumento apresentou crescimento de 52\% em quantidade de transações e respondeu por 47\% dos pagamentos sem uso de espécie no último trimestre do ano \cite{bcb_pagamentos_2024}. Esses números caracterizam escala e adoção, mas não devem ser interpretados como medida da incidência de fraude.

O risco associado a uma transferência não decorre apenas do funcionamento técnico do instrumento. Manipulação do usuário, comprometimento de credenciais, abuso de conta e mudança de comportamento podem produzir operações aparentemente regulares na camada transacional. Consequentemente, um classificador não observa diretamente a intenção do agente nem substitui a apuração de uma instituição financeira. Sua saída deve ser apresentada como estimativa de risco sujeita a revisão.

O Mecanismo Especial de Devolução foi introduzido no Regulamento do Pix pela Resolução BCB nº 103/2021 \cite{bcb_resolucao103_2021}. O mecanismo compõe um processo regulatório e operacional; não equivale a um rótulo de aprendizado de máquina, a uma garantia automática de devolução ou a uma conclusão jurídica de fraude. Como o Regulamento do Pix e o Guia do MED são atualizados, afirmações sobre procedimentos e prazos exigem versão e data de vigência \cite{bcb_resolucao1_2020,bcb_guia_med_v43}.

A Pesquisa Febraban de Tecnologia Bancária 2024 registra a presença de detecção de fraude e lavagem de dinheiro entre casos de uso de inteligência artificial declarados pelos bancos participantes \cite{febraban2024}. Trata-se de evidência setorial de pertinência, não de uma taxa de fraude Pix nem de validação de um algoritmo específico. Essa distinção é necessária para evitar que estatísticas de investimento ou adoção sejam transformadas em evidência de desempenho.

\section{Aprendizado de máquina para detecção de fraude}

Detecção de fraude financeira é um problema de classificação rara, dinâmica e sensível ao protocolo de avaliação. Entre as dificuldades recorrentes estão desbalanceamento, grande volume, adaptação dos fraudadores, atraso na confirmação do rótulo e custos assimétricos de erro \cite{abdallah2016survey,dalpozzolo2018realistic}. A ordem e o histórico das transações também carregam informação, o que motiva atributos de recência, frequência e desvio calculados somente com eventos anteriores \cite{jurgovsky2018sequence}.

Em conjuntos fortemente desbalanceados, a acurácia pode permanecer alta mesmo quando o classificador falha na classe positiva. Curvas e áreas de precisão--revocação são particularmente informativas nesse cenário \cite{saito2015precision}. Por isso, este trabalho prevê AUC-PR, AUC-ROC, precisão, revocação, F1 e matriz de confusão, acompanhadas do limiar usado. O limiar não deve ser escolhido no conjunto de teste.

A regressão logística será usada como linha de base por oferecer um ponto de comparação simples e reproduzível. Random Forest representa uma alternativa não linear baseada na agregação de árvores \cite{breiman2001random}, enquanto XGBoost usa boosting de árvores com mecanismos de regularização e escalabilidade \cite{chen2016xgboost}. Comparar esses modelos permite verificar se ganhos de complexidade se traduzem em melhoria relevante para a classe rara.

SMOTE cria amostras sintéticas da classe minoritária a partir de vizinhos no espaço de atributos \cite{chawla2002smote}. A técnica não deve ser aplicada antes da divisão dos dados: nesse caso, exemplos derivados da validação ou do teste podem contaminar o treino. A comparação adequada coloca reamostragem dentro do pipeline ajustado somente em treino e a contrasta com alternativas como ponderação de classe. Variáveis categóricas, valores ausentes e atributos de alta cardinalidade exigem tratamento compatível antes de qualquer interpolação.

Mudança de conceito limita a interpretação de uma divisão aleatória como evidência operacional. Um modelo pode reconhecer padrões presentes simultaneamente em treino e teste sem manter desempenho em período posterior. Assim, uma separação temporal é metodologicamente preferível quando a ordem relativa está disponível, ainda que \texttt{TransactionDT} não revele a data civil real. O IEEE-CIS será usado como proxy técnico de comércio eletrônico/cartão \cite{ieee_cis_2019}; seus resultados não demonstram generalização para Pix.

\section{Explicabilidade por SHAP}

SHAP expressa uma previsão como valor de referência acrescido das contribuições atribuídas aos atributos \cite{lundberg2017unified}. A análise local ajuda a identificar quais entradas elevaram ou reduziram o escore em uma transação, enquanto agregações de valores absolutos apoiam uma visão global do comportamento do modelo.

Essas contribuições não estabelecem causalidade, intenção criminosa ou significado semântico de uma variável anônima. No IEEE-CIS, grupos como \texttt{C*}, \texttt{D*}, \texttt{M*}, \texttt{V*} e \texttt{id\_*} possuem semântica individual parcial ou não publicada. Uma explicação correta pode informar que \texttt{V258} influenciou o modelo, mas não pode renomeá-la como conta, dispositivo ou destinatário Pix sem uma fonte que sustente essa equivalência.

\section{Retrieval-Augmented Generation e modelos de linguagem}

Modelos de linguagem organizam texto a partir de padrões aprendidos, mas não garantem atualização, rastreabilidade ou fidelidade factual. Retrieval-Augmented Generation combina recuperação de documentos com geração condicionada ao contexto recuperado \cite{lewis2020rag}. Em vez de depender apenas da memória paramétrica do modelo, o sistema consulta uma coleção externa e inclui trechos relevantes na entrada da geração.

Um pipeline de recuperação textual contém documentos, segmentação em trechos, embeddings, índice e retriever. Sentence-BERT mostrou uma forma eficiente de produzir representações de sentenças adequadas à comparação semântica \cite{reimers2019sentencebert}, enquanto recuperação densa representa consultas e passagens em espaços vetoriais compatíveis \cite{karpukhin2020dense}. FAISS oferece algoritmos e estruturas para busca eficiente por similaridade vetorial \cite{johnson2021billion}.

O retriever recebe uma consulta e devolve trechos candidatos; o modelo de linguagem recebe esses trechos e produz a resposta. O banco vetorial, contudo, não substitui o documento original como fonte de verdade. Cada trecho precisa preservar identificador, emissor, versão, vigência, página ou seção e URL. No domínio regulatório, misturar versões históricas e vigentes pode produzir uma resposta linguisticamente plausível e operacionalmente incorreta.

RAG também não elimina alucinação. A recuperação pode retornar material irrelevante, e o gerador pode produzir afirmações não sustentadas pelo contexto. A avaliação deve separar qualidade da recuperação, fidelidade ao documento, coerência com a evidência SHAP e clareza humana. Métricas automatizadas, como as propostas pelo RAGAS, podem apoiar essa análise, mas não substituem um conjunto de referência e revisão humana \cite{es2024ragas}.

\section{Síntese e lacuna investigada}

A literatura fornece métodos para classificação rara, explicação local e recuperação documental, porém sua simples justaposição não garante uma explicação fiel. O classificador responde quanto o padrão se aproxima da classe aprendida; SHAP registra quais entradas influenciaram a previsão; o RAG recupera contexto verificável. A lacuna investigada é integrar essas evidências sem transformar correlação em causalidade, sem inventar semântica para atributos anonimizados e sem confundir risco estatístico com decisão regulatória.

O protótipo adotará um registro versionado de atributos explicáveis. Apenas features explícitas ou derivadas com fórmula, janela, unidade e limitação aprovadas poderão fornecer conceitos ao retriever. Colunas anônimas podem permanecer úteis ao classificador, mas serão descritas como atributos anonimizados quando aparecerem na explicação.

\chapter{Metodologia}

\section{Desenho da pesquisa}

O estudo foi conduzido como prova de conceito experimental de um pipeline híbrido para classificação e explicação de risco. O conjunto IEEE-CIS Fraud Detection foi usado para desenvolvimento e avaliação técnica \cite{ieee_cis_2019}. Como os dados representam comércio eletrônico/cartão e contêm atributos anonimizados, o experimento não permite estimar desempenho operacional em transações Pix reais. A transferência ao contexto Pix foi limitada a papéis analíticos explicitamente documentados.

O fluxo possui três camadas: classificação por aprendizado de máquina, explicação local por SHAP e contextualização documental por RAG. Cada camada foi definida com entrada, saída, versão e protocolo de avaliação próprios. Essa separação evita que uma narrativa produzida pelo modelo de linguagem seja confundida com a previsão do classificador ou com o conteúdo de uma norma.

\section{Preparação dos dados e protocolo experimental}

As tabelas de transação e identidade foram unidas por \texttt{TransactionID} por meio de junção à esquerda, preservando todas as transações. A análise local encontrou 590.540 transações, das quais 20.663 apresentam \texttt{isFraud=1}, correspondendo a 3,499\%. A tabela de identidade possui 144.233 linhas e cobre 24,424\% das transações do treino.

As observações foram ordenadas por \texttt{TransactionDT}, que representa um delta temporal em segundos a partir de uma referência não publicada. O valor preserva ordem e intervalos relativos, mas não foi convertido em data civil ou horário local. A avaliação principal adotou uma divisão temporal de 70\% para treino, 15\% para validação e 15\% para teste, com as transações mais antigas no treino e as mais recentes no teste. Transações com o mesmo valor de \texttt{TransactionDT} permaneceram na mesma partição. Transformações, imputação, \emph{encoding}, escalonamento, seleção e estatísticas históricas foram ajustados somente no treino.

No conjunto completo, o corte resultou em 413.378 transações de treino, 88.581 de validação e 88.581 de teste, contendo, respectivamente, 14.538, 3.042 e 3.083 registros marcados como fraude. As taxas correspondentes foram 3,517\%, 3,434\% e 3,480\%, uma diferença máxima de 0,083 ponto percentual. Portanto, a classe positiva permaneceu rara, mas sua proporção foi semelhante entre as três partições sem que o alvo fosse usado para estratificação. Ocorrências repetidas de \texttt{TransactionDT}, desconsiderando a primeira de cada valor, corresponderam a 2,9\% das linhas; ao contar todas as linhas pertencentes a grupos de \emph{timestamps} repetidos, a proporção foi 5,7\%. Nenhuma ocorrência empatada atravessou as fronteiras do corte. Uma análise aleatória estratificada pode ser apresentada apenas como complemento, sem substituir a avaliação temporal principal.

O registro \texttt{config/pix\_feature\_registry.json} documenta cada atributo derivado: fórmula, colunas-fonte, população, janela, tratamento de ausentes, controle de vazamento, tags semânticas e advertência. Em decisão registrada pela dupla em 16/08/2026, foram aprovados e implementados quatro conceitos: desvio robusto do valor no histórico proxy, frequência recente no identificador proxy, raridade de dispositivo nesse histórico e posição cíclica temporal relativa. \texttt{card1} é tratado como proxy mascarado do domínio original, nunca como conta ou chave Pix; a origem e o fuso de \texttt{TransactionDT} não são presumidos.

A regressão logística foi usada como linha de base. Ponderação de classe e SMOTE foram comparados em configurações separadas, com a reamostragem restrita à partição de treino \cite{chawla2002smote}. Para o artefato XGBoost final, o limiar foi selecionado na validação pela maximização do F1 e congelado antes da avaliação do teste. Foram registradas AUC-PR, AUC-ROC, precisão, revocação, F1 e matriz de confusão no teste preservado \cite{saito2015precision}.

\section{Camada 1: classificação de risco}

A primeira camada recebe os atributos preparados e produz a probabilidade e a classe de risco. XGBoost foi adotado como modelo principal, enquanto Random Forest e regressão logística foram usados como comparadores \cite{breiman2001random,chen2016xgboost}. O artefato XGBoost é distribuído com o pré-processador e o manifesto, e sua política de decisão foi persistida separadamente. O retorno atual do pipeline inclui classe, probabilidade e limiar, mas não expõe no mesmo pacote a versão do modelo e a versão do registro de atributos.

\section{Camada 2: explicação local por SHAP}

A segunda camada calculou contribuições SHAP para o XGBoost em transações selecionadas \cite{lundberg2017unified}. O método foi escolhido por fornecer atribuições locais aditivas e uma agregação global pela média dos valores absolutos. Depois de congelados o conjunto de dados, o pré-processamento, a ordem das colunas e o modelo, o \texttt{TreeExplainer} foi configurado com \texttt{feature\_perturbation=interventional}, fundo de 500 linhas amostradas somente do treino por \emph{seed} registrada e \texttt{model\_output=probability}. Essa escala permite verificar a soma das contribuições contra a probabilidade prevista; a saída bruta padrão do XGBoost estaria em margem ou \emph{log-odds}. A versão do SHAP foi registrada explicitamente. O Random Forest não recebeu, nesta etapa, uma análise SHAP equivalente e não foi incluído na comparação de explicações locais.

Para cada observação, foi verificada a propriedade aditiva
\begin{equation}
    f(x) \approx E[f(X)] + \sum_{j=1}^{p} \phi_j,
\end{equation}
em que $E[f(X)]$ é o valor-base e $\phi_j$ representa a contribuição do atributo $j$ na mesma escala de $f(x)$. O erro máximo de reconstrução registrado foi aproximadamente $3{,}04 \times 10^{-7}$, dentro da tolerância numérica adotada. Valores não finitos ou desalinhamento entre nomes e posições de colunas invalidam o artefato explicativo.

A análise global ordenou os atributos pela média de $|\phi_j|$ em 1.000 observações da validação. A análise local registrou casos representativos de verdadeiro positivo, falso positivo, falso negativo e verdadeiro negativo segundo o limiar histórico de 0,5. A seleção desses casos não foi refeita com o limiar final de 0,6834 e, por isso, seus rótulos locais não devem ser reinterpretados no novo ponto de operação. As contribuições são tratadas como influência estatística no comportamento do modelo, não como causalidade, culpa ou intenção.

Para conectar SHAP ao RAG, o pipeline seleciona os três fatores de maior contribuição absoluta, preservando nome, valor transformado, contribuição e direção. O desenho metodológico previa que somente atributos com semântica aprovada gerassem \emph{tags} conceituais. A implementação atual, contudo, constrói a consulta a partir dos nomes dos fatores, inclusive nomes anônimos, sem lhes atribuir significado Pix. Essa diferença foi mantida explícita para que uma limitação do protótipo não seja descrita como salvaguarda já implementada.

A avaliação verificou fidelidade numérica e alinhamento entre os atributos e suas posições. A consistência entre os valores SHAP e a narrativa posterior permanece sem resultado, pois o cliente LLM não foi implementado. A distribuição das contribuições pode depender da configuração do explicador, do conjunto de referência e da correlação entre atributos. Uma explicação fiel ao modelo não garante que o modelo esteja correto, calibrado ou livre de viés; por isso, SHAP é tratado como instrumento de auditoria do classificador e não como validação independente da decisão.

\section{Camada 3 --- Retrieval-Augmented Generation (RAG)}

A terceira camada prepara a contextualização documental da saída numérica das etapas anteriores. A arquitetura de \emph{Retrieval-Augmented Generation} combina a recuperação de trechos externos com a geração de linguagem \cite{lewis2020rag}. Neste trabalho, o RAG não altera a classe ou a probabilidade calculada pelo modelo, não substitui a explicação SHAP e não decide se houve fraude. Seu papel é recuperar contexto rastreável para a predição e para as contribuições locais, mantendo separadas a evidência do classificador, a informação documental e as limitações da análise.

A base de conhecimento é composta por quatro fontes catalogadas: o Regulamento do Pix estabelecido pela Resolução BCB nº 1/2020, a Resolução BCB nº 103/2021, o Guia do Mecanismo Especial de Devolução na versão 4.4 e a Pesquisa Febraban de Tecnologia Bancária 2024 --- volume 1 \cite{bcb_resolucao1_2020,bcb_resolucao103_2021,bcb_guia_med_v44,febraban2024}. As três primeiras fontes oferecem conteúdo normativo ou operacional, enquanto o relatório da Febraban é empregado apenas como contexto setorial. Para manter a rastreabilidade, cada documento conserva identificador, título, emissor, tipo, URL, data de publicação, versão, situação normativa, datas de vigência e hash SHA-256. Cada trecho derivado também preserva sua página ou seção e a referência ao documento de origem.

Na etapa de indexação, os arquivos PDF e JSON são extraídos e normalizados. O conteúdo é inicialmente segmentado em \emph{chunks} pais de 500 tokens, com sobreposição de 50 tokens. Como o modelo de \emph{embeddings} possui limite de 128 tokens, cada \emph{chunk} pai é dividido em janelas de até 126 tokens de conteúdo, com sobreposição de 24 tokens. As representações vetoriais são produzidas pelo modelo multilíngue \texttt{sentence-transformers/paraphrase-multilingual-MiniLM-L12-v2}, na revisão registrada pelo manifesto do índice, com vetores de 384 dimensões \cite{reimers2019sentencebert}. Após normalização L2, os vetores são armazenados em um índice FAISS exato do tipo \texttt{IndexFlatIP}; nessa configuração, o produto interno entre vetores normalizados equivale à similaridade de cosseno \cite{johnson2021billion}. A reconstrução versionada de 20 de setembro de 2026 reuniu 106 unidades extraídas, 290 \emph{chunks} pais e 1.195 janelas vetoriais. O índice, os metadados e o manifesto de integridade foram versionados. Três fontes mantiveram os mesmos hashes do registro histórico; o arquivo consolidado do Regulamento do Pix apresentou hash distinto e, por isso, a reconstrução não foi tratada como idêntica à base anterior.

A recuperação recebe uma consulta construída com a situação prevista e os nomes dos três fatores de maior contribuição absoluta identificados pelo SHAP. Nomes anônimos são preservados literalmente, sem tradução para conceitos de negócio. A consulta é codificada pelo mesmo modelo de \emph{embeddings} e comparada aos vetores do índice. O recuperador devolve os cinco trechos de maior similaridade, acompanhados de escore e metadados. A implementação não aplica limiar mínimo de similaridade nem filtro automático de vigência. Consequentemente, a presença de um trecho entre os cinco primeiros resultados não comprova sua relevância ou sua aplicabilidade temporal, e fontes com vigência escalonada exigem conferência humana antes de fundamentar uma afirmação.

A integração entre SHAP e RAG ocorre no orquestrador \texttt{src.pipeline.explicar\_transacao}. A classificação e a probabilidade são calculadas pelo XGBoost, e os três fatores de maior contribuição absoluta são encaminhados ao módulo de recuperação. O contexto do gerador preserva o nome, o sinal e o valor da contribuição SHAP, mas a versão atual do prompt não inclui o valor bruto observado da variável, a versão do modelo ou a vigência validada de cada documento. A integração, portanto, não autoriza causalidade e mantém a diferença entre o domínio IEEE-CIS e transações Pix reais, porém ainda necessita desses campos para cumprir integralmente o pacote de evidências projetado.

O prompt é montado por \texttt{src.rag.explainer.montar\_prompt}. Ele contém instruções de comportamento, a decisão e a probabilidade do classificador, os três fatores SHAP com a direção e a contribuição numérica e até cinco trechos recuperados, identificados pela fonte. As instruções exigem resposta em português do Brasil, com três a cinco frases, proíbem afirmações causais e a atribuição de significado às colunas anônimas, restringem a resposta às evidências fornecidas e determinam abstenção quando o suporte for insuficiente. Os metadados completos de URL, localização, versão, vigência e escore são mantidos no retorno estruturado do pipeline e podem ser exibidos na interface, mas não são todos incorporados ao texto atual do prompt.

O contrato do cliente de modelo de linguagem foi definido para receber o prompt e devolver texto, mas sua implementação permanece \texttt{\{\{PREENCHER\}\}}. A decisão registrada indica Claude Haiku 4.5 como opção principal e Ollama local como alternativa; nenhuma das duas integrações foi implementada ou avaliada. Na ausência do cliente, o pipeline não simula uma resposta: informa que a explicação em linguagem natural está indisponível, preserva o prompt montado para inspeção e mantém visíveis a predição, os fatores SHAP e os documentos recuperados.

As estratégias implementadas para reduzir alucinações são a fundamentação estrita, o limite do contexto aos trechos recuperados, a preservação das contribuições numéricas, a proibição de semântica inventada para variáveis anônimas e a regra de abstenção. O prompt também orienta que a transação pertence a um conjunto de pesquisa e que ninguém seja tratado como culpado. O controle automático de vigência e a validação factual pós-geração não foram implementados. A avaliação proposta considera relevância da recuperação, coerência com SHAP, fidelidade aos dados, ausência de alucinações e clareza em português \cite{es2024ragas}; entretanto, até o fechamento desta versão, somente consultas de fumaça haviam verificado o carregamento e o funcionamento do índice, sem produzir métricas de qualidade ou respostas de LLM.

\section{Avaliação integrada}

A classificação foi avaliada por métricas adequadas à classe rara, e a explicação SHAP foi verificada quanto à aditividade e ao alinhamento dos atributos. Para a recuperação, foi definido um protocolo com relevância dos trechos, precisão em $k$, \emph{recall} em $k$ e posição do primeiro resultado relevante. Para a geração, foram definidos critérios de suporte documental, coerência com SHAP, fidelidade à transação, clareza e ausência de linguagem acusatória. Como o conjunto anotado e o cliente LLM não estavam disponíveis no fechamento desta versão, a avaliação experimental de recuperação e geração permanece \texttt{\{\{RESULTADO\}\}}; consultas de fumaça não foram interpretadas como medida de qualidade. Avaliações automáticas podem apoiar, mas não substituir, a inspeção humana \cite{es2024ragas}.

\section{Reprodutibilidade, segurança e limitações}

O código, os parâmetros, as \emph{seeds}, as versões e as decisões metodológicas foram registrados no repositório. Os arquivos brutos do IEEE-CIS, a matriz de fundo do SHAP, os exemplos extraídos do conjunto e eventuais chaves de API permanecem fora do Git. Exemplos enviados a serviços externos devem ser sintéticos ou minimizados. Saídas do modelo de linguagem não devem executar comandos nem produzir decisão automática de bloqueio, acusação ou devolução.

As principais ameaças à validade são a diferença de domínio entre IEEE-CIS e Pix, a anonimização de atributos, a mudança de distribuição, a confirmação imperfeita do rótulo e a ausência de validação em uma instituição financeira. Somam-se a elas a falta de avaliação anotada do RAG, a ausência do cliente LLM e a diferença entre o limiar de 0,6834 selecionado na validação do artefato final e o valor padrão de 0,5 atualmente utilizado pela interface. Portanto, os resultados constituem evidência parcial de viabilidade técnica da integração e não caracterizam um sistema pronto para produção.

\chapter{Resultados e Discussão}

\texttt{\{\{REVISÃO\_NECESSÁRIA\}\}} O texto consolidado deste capítulo não foi disponibilizado para a revisão de setembro. Antes da entrega final, devem ser inseridos e discutidos os resultados já registrados nos artefatos do projeto: comparação entre os classificadores, avaliação do XGBoost no conjunto de teste preservado, análises SHAP global e local, consultas de recuperação documental e comportamento da interface. A seção não deve apresentar exemplos de geração por LLM enquanto o cliente não estiver implementado nem preencher métricas do RAG sem avaliação anotada.

\chapter{Conclusão}

\section{Síntese do trabalho}

Este trabalho investigou a integração de aprendizado de máquina, explicabilidade por SHAP e recuperação documental em uma prova de conceito voltada à explicação de risco de fraude. O desenvolvimento foi realizado sobre o conjunto IEEE-CIS Fraud Detection, cuja natureza de comércio eletrônico e cujos atributos parcialmente anonimizados impedem a extrapolação direta dos resultados para transações Pix. Assim, a contribuição deve ser interpretada como demonstração técnica e metodológica, não como validação de um sistema antifraude para uso financeiro real.

O pipeline de dados preservou a ordem definida por \texttt{TransactionDT} e separou treino, validação e teste na proporção aproximada de 70/15/15, mantendo valores temporais empatados na mesma partição. O XGBoost foi adotado como classificador principal e comparado com a regressão logística e o Random Forest. Para o artefato final versionado, o limiar de 0,6834 foi selecionado exclusivamente na validação pela maximização do F1. Nessa partição, a AUC-PR foi 0,5719. No teste temporal preservado, o modelo obteve AUC-PR de 0,4833, AUC-ROC de 0,8700, precisão de 0,5852, revocação de 0,4324 e F1 de 0,4973. A redução de 0,0887 na AUC-PR entre validação e teste evidencia degradação temporal no próprio conjunto estudado e reforça a necessidade de monitoramento e revalidação antes de qualquer uso externo.

A camada SHAP foi aplicada ao XGBoost com \texttt{TreeExplainer}, fundo de 500 observações do treino, perturbação interventional e saída na escala de probabilidade. A verificação de fidelidade apresentou erro máximo de reconstrução de aproximadamente $3{,}04 \times 10^{-7}$. Na análise global registrada, \texttt{TransactionAmt}, \texttt{C13}, \texttt{C1}, \texttt{C14} e \texttt{card1} ocuparam as cinco primeiras posições por média do valor SHAP absoluto. Esses resultados descrevem o comportamento do classificador; não atribuem causalidade nem autorizam inferir o significado de colunas anônimas.

A camada documental reuniu quatro fontes normativas, operacionais e setoriais. A reconstrução versionada produziu 290 \emph{chunks} pais e 1.195 janelas vetoriais de 384 dimensões em um índice FAISS exato. O recuperador, o construtor de prompt, o orquestrador do pipeline e a interface Streamlit foram implementados. A interface recebe uma transação, apresenta a predição, os principais fatores SHAP e os documentos recuperados e informa a indisponibilidade de componentes ausentes. O cliente do modelo de linguagem, contudo, permanece \texttt{\{\{PREENCHER\}\}}. Por essa razão, não foram produzidas explicações reais de LLM nem resultados experimentais de coerência, fidelidade, alucinação ou clareza. As consultas executadas sobre o índice foram testes de fumaça e não constituem avaliação de relevância.

\section{Atendimento aos objetivos}

Os objetivos de preparação e caracterização do IEEE-CIS, comparação de modelos, explicação por SHAP, registro de atributos e organização da base documental foram atendidos nos limites da prova de conceito. Também foram implementadas a recuperação semântica, a montagem controlada do prompt e a exposição das camadas na interface.

O objetivo relativo à geração de explicações foi atendido apenas quanto ao contrato do cliente, à preparação do contexto e às salvaguardas do prompt. A chamada a um LLM não foi implementada. De modo semelhante, foi definido um protocolo objetivo para avaliar recuperação e geração, mas sua execução permanece \texttt{\{\{RESULTADO\}\}} por falta do cliente LLM e de um conjunto anotado. O objetivo de reprodutibilidade foi parcialmente atendido: código, modelo principal e índice vetorial foram versionados, enquanto os dados brutos, os exemplos da interface e a matriz de fundo do SHAP permanecem externos por licença, tamanho ou método de geração documentado.

Também subsiste uma divergência operacional. O modelo final foi avaliado com limiar de 0,6834, mas o pipeline usado pela interface mantém 0,5 como valor padrão. Até que a aplicação carregue explicitamente o limiar persistido do artefato, suas classificações não devem ser apresentadas como reprodução do ponto de operação avaliado no teste.

\section{Contribuições}

A principal contribuição do trabalho é uma arquitetura modular na qual previsão, atribuição local e recuperação documental permanecem distinguíveis. Essa separação favorece a auditoria: a probabilidade e a classe pertencem ao classificador; as contribuições pertencem ao explicador; e as afirmações contextuais devem ser sustentadas pelas fontes recuperadas. O projeto também contribui ao preservar a ausência de significado das variáveis anonimizadas, tratar a versão dos documentos como parte da evidência e declarar indisponibilidade em vez de gerar uma saída simulada.

No plano experimental, o uso de divisão temporal, métricas adequadas à classe rara e escolha do limiar somente na validação permitiu observar perda de desempenho em dados posteriores. No plano de reprodutibilidade, a disponibilização do artefato XGBoost e do índice FAISS reduz a dependência de retreinamentos e reconstruções potencialmente divergentes.

\section{Limitações}

As limitações centrais são a diferença entre os domínios IEEE-CIS e Pix, a anonimização das variáveis, a ausência de dados transacionais reais de instituições financeiras e a falta de validação por especialistas do domínio. O rótulo do conjunto não elimina problemas de atraso ou imperfeição de confirmação, e a queda de desempenho no teste indica sensibilidade à mudança temporal.

Na camada explicativa, SHAP descreve o modelo, não a realidade causal da fraude. A matriz de fundo necessária à configuração interventional não é versionada e precisa ser reproduzida pelo script indicado no projeto. Na camada RAG, a consulta atual usa os nomes dos fatores mais influentes, inclusive nomes anônimos, sem aplicar limiar de similaridade ou filtro automático de vigência. Um dos documentos da reconstrução atual também possui hash diferente do registro histórico. Por fim, a ausência do cliente LLM impede verificar empiricamente se as restrições do prompt são suficientes para evitar alucinações ou contradições.

\section{Trabalhos futuros}

Como continuação imediata, recomenda-se carregar no pipeline o limiar de 0,6834 associado ao artefato avaliado e registrar essa informação em seu manifesto. Em seguida, deve-se implementar o cliente de LLM escolhido, com credenciais externas ao código, limites de custo, registro de versão e tratamento de indisponibilidade. A integração deve encaminhar ao prompt os valores observados, as versões do modelo e do registro de atributos e os metadados completos de cada fonte.

Também é necessário construir um conjunto de consultas e trechos relevantes anotados, executar o protocolo de recuperação e submeter as explicações a avaliação humana quanto à coerência com SHAP, fidelidade à transação, suporte documental, ausência de alucinações e clareza em português. Recomenda-se acrescentar filtros de vigência, critérios de abstenção baseados na qualidade da recuperação e testes adversariais contra instruções maliciosas presentes nos documentos.

Por fim, a evolução para um cenário aplicado exige dados Pix devidamente governados, revisão jurídica e de privacidade, validação com analistas de fraude, calibração do modelo, definição de custos de erro e monitoramento contínuo de desempenho e mudança de distribuição. Essas etapas pertencem a trabalhos futuros e não devem ser interpretadas como funcionalidades já entregues pelo protótipo.


\postextual
\bibliography{referencias}

\end{document}
